% ============================================================================
% Résumés : français · anglais
% ============================================================================

% ---------- Résumé français ----------
\chapter*{Résumé}
\addcontentsline{toc}{chapter}{Résumé}
\markboth{Résumé}{Résumé}

\begin{infobox}[Résumé --- Français]
La supervision des infrastructures informatiques constitue aujourd'hui un
enjeu stratégique pour toute organisation exploitant des serveurs. Les outils
traditionnels de surveillance tels que Prometheus, Grafana ou Datadog se
concentrent sur l'affichage de métriques, laissant aux ingénieurs la charge
d'interpréter les données et de diagnostiquer manuellement les incidents.
Cette approche génère une \textit{fatigue d'alertes} importante et allonge
considérablement le temps moyen de résolution (\textsc{mttr}).

Ce rapport présente \textbf{\caimp{}} (\textit{Centralized AI Monitoring
Platform}), une plateforme open source auto-hébergée qui repositionne la
supervision autour d'une philosophie \textit{réponses-en-premier} : au lieu
d'afficher des graphes, le tableau de bord répond directement à la question
\og{}Quelque chose est-il cassé, et si oui, pourquoi ?\fg{}

L'architecture repose sur 19 microservices conteneurisés communiquant via
\textsc{nats} JetStream. Un agent léger déployé sur chaque serveur cible
collecte les métriques système via \textsc{otlp} et signale automatiquement
les changements d'environnement (déploiements Git, événements Docker,
installations de paquets, connexions \textsc{ssh}). Un moteur de corrélation
associe chaque anomalie détectée aux modifications survenues dans les 30
minutes précédant l'incident et soumet l'analyse à un \textsc{llm} local
(Llama 3.1 8B via Ollama) pour produire une explication en langage naturel.

Le moteur de prévision applique le lissage exponentiel de Holt-Winters sur
les 24 heures à venir. Un moteur de déduplication d'alertes élimine les
doublons sur une fenêtre de 15 minutes. La plateforme est installable en
moins de 10 minutes via Docker Compose.
\end{infobox}

\medskip
\textbf{Mots-clés :} AIOps, supervision d'infrastructure, microservices,
grand modèle de langage, TimescaleDB, NATS JetStream, corrélation de
changements, détection d'anomalies, Holt-Winters, RAG, pgvector.

\vspace{1.5cm}

% ---------- Abstract (English) ----------
\chapter*{Abstract}
\addcontentsline{toc}{chapter}{Abstract (English)}
\markboth{Abstract}{Abstract}

\begin{infobox}[Abstract --- English]
Infrastructure monitoring is a strategic challenge for any organisation
operating servers. Traditional tools such as Prometheus, Grafana or Datadog
focus on displaying metrics, leaving engineers to manually diagnose incidents.
This generates significant \textit{alert fatigue} and increases the mean time
to resolution (\textsc{mttr}).

This report presents \textbf{\caimp{}} (\textit{Centralized AI Monitoring
Platform}), a self-hosted open-source platform built around an
\textit{answers-first} philosophy: instead of graphs, the dashboard answers
directly ``\textit{Is anything broken, and if so, why?}''

The architecture uses 19 containerised microservices communicating via
\textsc{nats} JetStream. A lightweight agent collects metrics via \textsc{otlp}
and automatically reports system changes. A change correlation engine scores
each anomaly against changes in the 30-minute window before it, then calls a
local \textsc{llm} (Llama 3.1 8B via Ollama) for a plain-English explanation
and recommended action.

A Holt-Winters forecasting engine predicts resource saturation over 24 hours.
An alert deduplication engine eliminates duplicates in a 15-minute window.
The platform installs in under 10 minutes via Docker Compose.
\end{infobox}

\medskip
\textbf{Keywords:} AIOps, infrastructure monitoring, microservices, large
language model, TimescaleDB, NATS JetStream, change correlation, anomaly
detection, Holt-Winters, RAG, pgvector.
